% =============================================================================
%  Section 3 — System Architecture
% =============================================================================
\section{Neuron System Architecture}
\label{sec:arch}

This section presents the architecture of Neuron, a multi-tenant LLM agent
governance runtime designed around the principle of \emph{coupled
resource--security governance}. The architecture comprises four design
layers---a unified process abstraction (Section~\ref{sec:arch-abstraction}),
resource cgroups (Section~\ref{sec:arch-cgroup}), execution sandboxes
(Section~\ref{sec:arch-sandbox}), and a unified LLM gateway
(Section~\ref{sec:arch-gateway})---bound together by the coupled security
enforcement loop (Section~\ref{sec:arch-coupled}), which is the central
contribution of this work. We adopt the terminology of classical operating
systems---process, syscall, cgroup, OOM killer, seccomp, capability---to make
the analogies precise and the overhead analysis tractable. As stated in
Section~\ref{sec:intro}, this is a \emph{functional} analogy: the trust
boundary is the JVM/container, not a hardware ring. Neuron runs in a single
JVM process, and its governance state (\texttt{PermissionProfile},
\texttt{AuditChain}, \texttt{SpawnPrivilegeContext}) shares the same address
space as agent code, with \texttt{InheritableThreadLocal} providing
language-level isolation rather than hardware-enforced ring separation. The
only hardware-enforced boundary is the \texttt{DockerSandboxProvider} path
(Section~\ref{sec:arch-sandbox}), which places destructive tool execution in
a disposable container with dropped capabilities and a seccomp profile; this
path is probed against seven known escape techniques in F8
(Section~\ref{sec:lim-failure}). Defending against JVM escape itself is out
of scope and is the responsibility of the JVM and container runtimes; the
threat model (Section~\ref{sec:threat}) assumes the JVM boundary holds, as
any in-process governance design must.

% -----------------------------------------------------------------------------
\subsection{Unified Abstraction: The Agent as an OS Process}
\label{sec:arch-abstraction}

The foundation of Neuron is a single abstraction that recurses through every
subsystem: an \emph{agent is an OS process}. A conventional process binds a
program counter, a file-descriptor table, and a scheduling priority to an
execution context; a Neuron agent process binds a \emph{turn}---an atomic
interaction comprising a prompt, a sequence of tool invocations, and an
LLM-mediated control flow---to a runtime-managed execution context. Within this
abstraction, the LLM token plays the role of the CPU cycle: it is the
indivisible unit of resource consumption that the scheduler accounts for, that
the cgroup caps, and that the OOM killer defends against.

\paragraph{Inter-process communication.}
Agents never invoke one another directly. Instead, all cross-agent
communication is mediated by a runtime-level \emph{EventBus}, a publish/subscribe
substrate analogous to Linux's netlink and udev event mechanisms. The EventBus
serves two classes of consumers: external SSE clients attached through a runtime
stream endpoint, and internal subscribers registered as
\texttt{Consumer<String>} callbacks. Internal subscriber dispatch is
performed on \emph{virtual threads} spawned per handler, so a slow subscriber
cannot block the publisher or starve co-resident subscribers---the
fan-out is non-blocking by construction. This design choice has a direct
governance consequence: an unchecked broadcast can trigger unbounded virtual
thread fan-out, a resource pressure vector that motivates the rate-limiting
interceptor introduced in Section~\ref{sec:arch-coupled}.

\paragraph{Caller identity as a runtime primitive.}
Every in-runtime resource decision must answer ``who is issuing this call?''
Neuron answers this through \texttt{CallerContext}, a thread-local that mirrors
Linux's \texttt{current} macro. It carries the caller's agent identifier and
tenant identifier, and is implemented as an
\texttt{InheritableThreadLocal} so that virtual threads---including those
spawned by \texttt{AgentTool}, Neuron's \texttt{fork()+exec()} equivalent for
sub-agent creation---automatically inherit the parent's identity. The
\texttt{CallerContext} is set by the \texttt{QueryEngine} (the turn scheduler)
immediately before tool execution and cleared in a \texttt{finally} block to
prevent thread-pool leakage. Runtime daemons that do not set a
\texttt{CallerContext} are exempt from per-tenant accounting, mirroring the
privileged-agent exemption in classical syscall filters.

\paragraph{Dispatch modes.}
The EventBus exposes two complementary dispatch modes. The broadcast mode
treats all subscribers as equal and concurrent. A second
\emph{responsibility-chain} dispatcher orders handlers by priority and permits
a high-priority handler to \emph{accept} (short-circuit) the chain, which is
the mode used for security review and permission adjudication where ordered,
potentially aborting decisions are required. The two modes thus separate
\emph{notification} (best-effort, parallel) from \emph{enforcement} (ordered,
fail-closed), a separation that the coupled enforcement loop of
Section~\ref{sec:arch-coupled} exploits.

% -----------------------------------------------------------------------------
\subsection{Resource Cgroups: Token Quotas and the OOM Killer}
\label{sec:arch-cgroup}

Neuron treats the LLM token as the analogue of memory, and governs it with a
hierarchical cgroup controller realized by \texttt{CgroupManager} and
\texttt{CgroupNode}. At initialization, the manager constructs a default quota
tree rooted at \texttt{neuron-root} (1{,}000{,}000 tokens), with children
\texttt{agents} (500{,}000), \texttt{system} (200{,}000), and
\texttt{tools} (300{,}000). Each tenant agent receives a per-instance cgroup
defaulting to 50{,}000 tokens, parented under \texttt{agents}. This hierarchy
is the direct analogue of Linux's \texttt{memory.max} cgroup tree.

\paragraph{Soft and hard limits.}
Each \texttt{CgroupNode} enforces two thresholds. The \emph{hard limit}
(\texttt{tokenQuota}, the analogue of
\texttt{memory.limit\_in\_bytes}) is a ceiling whose violation is fatal: when
\texttt{consumeTokens} detects that consumption plus the requested amount would
exceed the quota, it raises a \texttt{TokenOomException}---a \emph{token
OOM}---which the manager translates into an \texttt{oomKill} of the offending
agent. The
\emph{soft limit} (default 80\% of the quota, the analogue of
\texttt{memory.soft\_limit\_in\_bytes}) is non-fatal: its violation raises a
\texttt{TokenSoftOomException} that triggers \emph{semantic folding} via the
\texttt{SomWindowController}, a reclamation mechanism that compacts the agent's
context to reclaim tokens rather than killing the process. This two-tier design
provides \emph{fault tolerance} under transient pressure---the runtime first
attempts graceful reclamation and resorts to termination only when reclamation
cannot restore headroom.

\paragraph{Transactional propagation.}
Consumption propagates up the tree in a two-phase-commit style: a child records
its consumption locally, then delegates to its parent; if the parent raises an
OOM condition, the child \emph{rolls back} its local accounting before
re-propagating the exception. This prevents the partial-bookkeeping races that
would otherwise allow a child to believe it has spent tokens that its parent
never granted. A dedicated \texttt{preCheckAndReserve} gateway integrates this
check into the syscall dispatch path \emph{before} the LLM is invoked, so that
an over-quota agent is denied the call rather than permitted to consume and
killed after the fact---an enforcement ordering that reduces wasted
expenditure.

\paragraph{Privilege and control.}
Two control knobs mirror Linux's \texttt{oom\_control} and scheduling
privileges. \texttt{disableOomKill} exempts a critical agent (for example, the
init daemon) from termination, degrading its OOM to a logged event. Conversely,
agents running at \texttt{REALTIME} priority bypass all cgroup limits, the
analogue of a runtime-privileged path. Every cgroup decision---hard kill, soft
reclamation, consumption---is written to the unified audit chain described in
Section~\ref{sec:arch-coupled} under the \texttt{CGROUP} layer tag, so that resource
events are causally linkable to the security events that accompany them.

% -----------------------------------------------------------------------------
\subsection{Execution Sandbox: WASM and Container Isolation}
\label{sec:arch-sandbox}

Tool execution in Neuron is confined by a \texttt{SandboxProvider} abstraction
that realizes the Ring~3 user-mode execution environment. Two providers
implement the interface, offering complementary isolation granularity.

\paragraph{In-process WASM isolation.}
\texttt{GraalWasmSandbox} executes compiled WebAssembly within the host JVM
using the GraalVM polyglot engine, with \texttt{wasm.MaxMemoryPages} capped at
1024 pages to impose a hard memory ceiling. Sandboxed code cannot reach host
resources directly; instead, it issues \emph{semantic syscalls}
(\texttt{\_\_neuron\_syscall} imports) that are bound to runtime-mediated host
functions, so every privileged operation traverses the
Ring~3$\rightarrow$Ring~0 boundary and is subject to the permission
interceptor. Execution is wrapped in a fault-tolerance envelope:
signal interception drains \texttt{SIGTERM}/\texttt{SIGINT} before entry, and
any \texttt{PolyglotException} or uncaught exception invokes
\texttt{handleSandboxFault}, which destroys the sandbox instance and raises a
\texttt{SIGSEGV}-equivalent fault to the agent. This \emph{fault-and-isolate}
discipline ensures that a misbehaving computation is contained without
corrupting the host, at the overhead cost of one sandbox instance per
execution.

\paragraph{Out-of-process container isolation.}
\texttt{DockerSandboxProvider} executes untrusted code in a disposable Docker
container for workloads requiring stronger isolation or a full language
runtime. Each invocation writes the script to a unique
\texttt{Files.createTempFile} path---avoiding the race conditions and
cross-invocation leakage that fixed paths would introduce---and mounts it
read-only into a container launched with a hardened parameter set: network
disabled (\texttt{--network none}, air-gapped), memory capped at 256\,MB
(\texttt{-m} with matching swap), CPU limited to half a core
(\texttt{--cpus=0.5}), process count capped (\texttt{--pids-limit 64},
defeating fork bombs), a read-only root filesystem (\texttt{--read-only}),
and a \texttt{tmpfs} \texttt{/tmp} mounted \texttt{noexec}. Three
security options provide defense in depth: \texttt{--cap-drop ALL} strips all
Linux capabilities, \texttt{--security-opt no-new-privileges} forbids
\texttt{setuid}-based escalation by child processes, and
\texttt{--security-opt seccomp=<profile>} applies a blacklist that blocks
container-escape syscalls such as \texttt{unshare}, \texttt{pivot\_root},
\texttt{mount}, and \texttt{ptrace}. The seccomp profile is lazily materialized
from a packaged resource to a host temp file under double-checked locking and
registered for cleanup on JVM exit. The provider thus composes
capability-drop, privilege-disabling, and syscall-filtering into a single
fail-closed execution boundary.

% -----------------------------------------------------------------------------
\subsection{Unified LLM Gateway and Adaptive Model Routing}
\label{sec:arch-gateway}

The governance layers above assume a single point at which LLM consumption
enters the runtime. Neuron consolidates this point as a \emph{unified LLM
gateway} that all tenants and frameworks traverse, eliminating the
per-tenant or per-framework SDK that layer-separated designs must otherwise
maintain. The gateway routes each request through an \emph{adaptive model
router}: a lightweight, low-latency model---cheap to invoke and suitable for
high-frequency calls---serves as the front-end classifier for the
\texttt{ReviewGate} and as a permission-auditing screener, while
compute-intensive generation is delegated to a frontier LLM. Routing is
decided by the system-overhead signal observed at the gateway, so that under
resource pressure a larger share of traffic is deflected to the lightweight
model, reducing marginal token expenditure precisely when the cgroup is near
its soft limit.

This end-to-end routing is more than a cost optimization: it is the
\emph{source-side} complement to the cgroup. Because every request traverses
a single gateway, the router enforces a \emph{central rate limit} on token
consumption before any generation begins, so that a runaway tenant is
throttled at the entry point rather than chased after the fact by the OOM
killer. The gateway's per-request rate limit and the cgroup's per-turn quota
thus operate as two scales of the same mechanism---one per request, one per
turn---both stamped with the per-turn trace identifier and both written to
the unified audit chain. This echoes the cgroup principle at the LLM
boundary: the token is governed as a first-class resource from the moment it
is requested, not only after it is spent.

% -----------------------------------------------------------------------------
\subsection{Coupled Security Enforcement under Resource Pressure}
\label{sec:arch-coupled}

We now present the central mechanism of this work. In a layer-separated
runtime, the resource, isolation, and permission layers each enforce their own
policy against disjoint audit trails; under resource pressure, the security
layer is either starved outright or loses the causal context needed to detect
that a privileged operation is being attempted inside a contention window
(Section~\ref{sec:bg-lim}). Neuron closes this gap with a \emph{coupled enforcement
loop} in which resource pressure does not merely trigger a resource response
but directly \emph{escalates the interception level} or \emph{circuit-breaks
\IO} at the security boundary, with every coupled decision sharing a single
trace identifier.

\paragraph{The audit substrate: a unified audit chain.}
The substrate that makes cross-layer correlation possible is the
\texttt{UnifiedAuditLog}, a fourth, cross-layer sink that aggregates
decisions from all governance layers---\texttt{CGROUP}, \texttt{SANDBOX},
\texttt{PERMISSION}, and a resource-rate-limiting layer
\texttt{RATELIMIT}---into a single JSONL ledger under one schema. Each entry
carries a \texttt{traceId} injected by \texttt{TraceContext} at the turn
boundary via an \texttt{InheritableThreadLocal} (with save/restore semantics
to support nested turns and prevent thread-pool leakage). Because the trace
identifier is inherited by the virtual threads that execute sub-agents and
subscriber callbacks, every governance decision across the entire turn---an
OOM kill, a sandbox fault, a permission denial, a rate-limit drop---lands in
the same causal frame. The ledger can be replayed by \texttt{traceId} to
reconstruct the cross-layer response timeline of a single attack.
\emph{We emphasize that the trace identifier provides audit correlation, not
online enforcement:} it makes the coupling \emph{auditable} (the causal chain
from resource pressure to security response can be reconstructed for any
turn), but it does not, by itself, change any permission verdict. The online
coupling---in which a resource-layer signal tightens the permission decision
in real time---is realized by the joint-decision mechanism described later in
this section. Recording is best-effort and \texttt{FileChannel}-appended
outside the syscall dispatcher, so the audit path never recurses into
permission checking and cannot itself become a starvation victim.

\paragraph{\IO{} circuit-breaking at the resource boundary.}
The first coupling response to pressure is to circuit-break \IO{} \emph{before}
it reaches the permission layer. Two token-bucket interceptors, modeled on
Linux traffic control and the runtime's \texttt{RateLimitSyscallFilter}, gate
the shared substrates at their entry points. \texttt{EventBusRateLimiter}
attaches to \texttt{EventBus.broadcast} and enforces a dual-dimension
bucket---per-agent (default 50/s, burst 60) and per-tenant---so that no single
agent or tenant can monopolize the SSE fan-out or the virtual-thread pool.
When the bucket is exhausted, the broadcast is silently dropped (broadcast is
fire-and-forget, so raising would block the caller) and the denial is
double-written to the audit chain and the semantic ETW. \texttt{VfsRateLimiter}
gates VFS operations with \emph{differentiated} limits: writes and mounts are
heavily throttled (20/s, burst 25) because they consume locks and disk \IO{},
while reads are lightly throttled (200/s, burst 240) because they take shared
locks. A write that exceeds its quota raises a \texttt{SecurityException} that
the agent observes as a hard failure, forcing it to back off or abandon the
operation. Both interceptors exempt \texttt{sys.*} channels, runtime daemons,
and privileged agents, confining throttling to tenant-driven \IO{}. The
consequence is that a resource-contention campaign is damped \emph{at the
source} rather than policed after damage---the residual pressure converges to
zero, as the evaluation in Section~\ref{sec:eval} demonstrates.

\paragraph{Auto-escalation of the interception level.}
The second coupling response is to \emph{raise} the security enforcement level
when pressure signals accumulate, rather than evaluating permission at a fixed
strictness. Three escalation pathways realize this.

First, the \texttt{ToolCircuitBreaker} monitors per-agent, per-tool failure
counts. When an agent registers five consecutive failures against the same
tool---a signature of either a runaway loop or an attacker probing a
misconfigured path---the breaker \emph{trips} and revokes that agent's access
to the tool for a 60-second cooldown, escalating the effective permission from
\emph{allow} to \emph{deny} without reconsulting the declarative policy. After
the cooldown the breaker enters a half-open state, admitting a single probe
request to test whether the underlying fault has cleared; success resets the
breaker, while failure re-trips it. The trip event is broadcast on the
\texttt{sys.telemetry.events} channel and surfaced through the hook system,
closing the loop back to the audit chain.

Second, the soft-OOM signal from the cgroup layer escalates resource handling
from \emph{consume} to \emph{reclaim}: rather than admitting further
consumption and risking a hard kill, the runtime triggers semantic folding to
recover tokens, and only re-attempts the reservation if reclamation succeeds.
This is fault-tolerant back-pressure---the resource layer communicates pressure
to the execution layer, which responds by reducing its own demand.

Third, and most specific to the LIM attack surface, a depth-aware
\texttt{EscalationPolicy} escalates the permission verdict from \emph{ask} to
\emph{auto-deny} for destructive tools originating from deeply nested
sub-agents\label{sec:arch-lim}. When the \texttt{QueryEngine} reaches an
\texttt{ASK} decision, it gathers the caller's position in the spawn tree
(depth and parent chain from the \texttt{DelegationGuard}) and consults the
\texttt{EscalationPolicy}: if the depth meets or exceeds the configured
threshold (default 2) and the requested tool is destructive
(\texttt{bash}, \texttt{shell}, \texttt{security\_scan}), the verdict becomes
\texttt{DENY\_DEPTH}---the request is denied without consulting a human
approver, defeating the social-engineering pathway in which a deeply nested
sub-agent phrases an escalation in natural language. The denial is recorded
through the triple-write path (semantic ETW, unified audit log, and permission
denial ledger) together with the spawn-tree context, so that the escalation
attempt is auditable end-to-end. For non-destructive requests, or for
top-level calls at depth zero, the policy degrades gracefully to
\texttt{ASK\_WITH\_CONTEXT}, preserving the existing approval contract and
ensuring zero regression for headless workloads.

\paragraph{Resource-pressure-aware joint decision.}
The static \texttt{EscalationPolicy} above is a pure function of
(depth, tool name) and does not consult resource state. To realize
\emph{online} cross-layer coupling---where a resource-layer signal
actually changes the permission verdict in real time---we extend the
policy with a joint-decision mode in
\texttt{ResourcePressureAwareEscalationPolicy}. When the resource
layer's cumulative rate-limit rejection count (maintained by the
\texttt{VfsRateLimiter} and \texttt{EventBusRateLimiter} on the
per-turn trace) exceeds a configured \texttt{PRESSURE\_THRESHOLD}
(default 50, indicating sustained attacker pressure), the permission
layer dynamically tightens its strictness from depth~$\geq 2$ to
depth~$\geq 1$ for destructive tools. Concretely, a depth~1 sub-agent
requesting \texttt{bash} receives \texttt{ASK\_WITH\_CONTEXT} under
normal resource conditions but \texttt{DENY\_DEPTH} once resource
pressure exceeds the threshold---the resource signal \emph{changes}
the permission verdict. This is the mechanism that distinguishes
\emph{audit correlation} (the trace identifier's role, which is
post-hoc) from \emph{online enforcement} (the joint decision's role,
which is real-time). The threshold and tightening step are
heuristically calibrated and their sensitivity is ablated in
Section~\ref{sec:eval-joint}; we do not claim the rule generalizes to
all attack morphologies (L7, Section~\ref{sec:conc-lim}).

\paragraph{Spawn-time privilege non-increase.}
Coupling is also enforced \emph{across the spawn boundary}, closing the
attack surface in which a downgraded parent produces an upgraded child. An
\texttt{SpawnPrivilegeContext}, itself an \texttt{InheritableThreadLocal}
carrying the parent's effective \texttt{PermissionProfile}, is published by the
\texttt{QueryEngine} before tool execution and inherited automatically by the
virtual thread that \texttt{AgentTool} spawns. The child \texttt{QueryEngine}
applies the inherited profile on construction, so that a parent restricted by a
\texttt{*:deny} rule cannot produce a child with default privileges: spawn is
guaranteed to be privilege-non-increasing, the LIM analogue of a
\texttt{NoNewPrivs} flag. Combined with tenant-identity propagation through
the \texttt{CallerContext} and a hard \texttt{checkTenantOwnership}
pre-check that cannot be bypassed by permissive modes, the spawn boundary
becomes a security-preserving, tenant-isolating interface rather than an
escalation opportunity.

\paragraph{Summary of the loop.}
In aggregate, the coupled enforcement loop operates as follows. A resource
pressure signal---soft OOM, rate-limit exhaustion, or a failure
cascade---is observed by the resource layer. Rather than treating it as a
local resource event, Neuron translates it into a security response:
\IO{} is circuit-broken at the substrate boundary, the interception level is
escalated from allow to deny (via the tool circuit breaker or the depth-aware
escalation policy), and the spawn boundary is guarded to prevent privilege
non-increase from being violated under the cover of contention. Every
decision along this loop shares the per-turn \texttt{traceId} and is written
to the \texttt{UnifiedAuditLog}, enabling the cross-layer causal
reconstruction that layer-separated designs cannot provide. The coupled
governance claim is validated in Section~\ref{sec:eval}.

% -----------------------------------------------------------------------------
\subsection{Invariant Argument for Spawn-Time Privilege Non-Increase}
\label{sec:arch-invariant}

The security property of Section~\ref{sec:arch-lim}---that a child agent's
privileges are always a subset of its parent's---is the LIM analogue of
the \texttt{NoNewPrivs} flag in Linux. We provide a non-formal but
rigorous argument that this invariant holds across all code paths,
including the exceptional and concurrent cases. The argument proceeds
by enumerating the operations that could potentially violate the
invariant, showing that each is structurally blocked.

\paragraph{Notation and invariant.}
Let $\mathit{prof}(a)$ denote the effective \texttt{PermissionProfile}
of agent $a$ at the moment it begins executing a turn, and let
$\mathit{prof}_{\mathrm{default}}$ denote the permissive default
profile assigned by \texttt{PermissionProfile.empty()}. The
\emph{spawn-time privilege non-increase} invariant is:
\[
  \forall a_{\mathrm{child}}.\ \mathit{prof}(a_{\mathrm{child}})
  \sqsubseteq \mathit{prof}(a_{\mathrm{parent}}),
\]
where $\sqsubseteq$ is the profile subsumption order (a profile $p$
subsumes $q$ iff every permission granted by $q$ is also granted by
$p$). The complementary failure mode is that the child receives
$\mathit{prof}_{\mathrm{default}}$ when the parent is restricted,
which would violate the invariant because
$\mathit{prof}_{\mathrm{default}} \not\sqsubseteq
\mathit{prof}_{\mathrm{parent}}$ for any restricted parent.

\paragraph{Code-path enumeration.}
There are exactly four ways a child \texttt{QueryEngine} could
obtain its initial profile; we show that none can violate the
invariant.

\begin{enumerate}
    \item \textbf{Inherit from \texttt{SpawnPrivilegeContext}.} The
    \texttt{AgentTool} spawns the child on a virtual thread; the
    child's \texttt{QueryEngine} constructor reads
    \texttt{SpawnPrivilegeContext.current()}, which was
    \texttt{set} by the parent's \texttt{QueryEngine} before
    tool execution and is inherited automatically by the
    virtual thread. The child thus receives
    $\mathit{prof}(a_{\mathrm{parent}})$, satisfying the invariant
    with equality.

    \item \textbf{Fall back to default when context is \texttt{null}.}
    If \texttt{SpawnPrivilegeContext.current()} returns \texttt{null},
    the child falls back to $\mathit{prof}_{\mathrm{default}}$.
    This path is reachable \emph{only} when no parent has set the
    context, i.e.\ for runtime daemons and headless tests; it is
    unreachable for tenant-driven spawns because the parent
    \texttt{QueryEngine} always sets the context before invoking
    \texttt{AgentTool}.

    \item \textbf{Exception during spawn.} If the spawn raises an
    exception after the context is set but before the child
    constructor reads it, the child is never constructed and the
    invariant holds vacuously. If the exception occurs after the
    child reads the context, the child aborts with the inherited
    profile, again satisfying the invariant.

    \item \textbf{Concurrent spawn from multiple parents.}
    \texttt{InheritableThreadLocal} provides per-thread isolation,
    so two parents spawning children concurrently do not interfere:
    each child inherits its own parent's profile. There is no shared
    mutable state on the spawn path.
\end{enumerate}

\paragraph{Pseudocode of the invariant-preserving spawn path.}
The structural argument above is realized by the following pseudocode,
which mirrors the Java implementation of \texttt{QueryEngine} and
\texttt{AgentTool}:

\begin{verbatim}
// Parent QueryEngine, before tool execution:
function executeTurn(turn):
    SpawnPrivilegeContext.set(self.effectiveProfile)
    try:
        result = turn.run()
    finally:
        SpawnPrivilegeContext.clear()   // force-clear on all paths
    return result

// AgentTool, spawning a child:
function spawnChild(prompt):
    // SpawnPrivilegeContext.current() is inherited by the
    // virtual thread that runs the child QueryEngine.
    childThread = Thread.startVirtualThread(() -> {
        child = new QueryEngine(prompt)
        child.run()
    })
    return childThread

// Child QueryEngine constructor:
constructor(prompt):
    inherited = SpawnPrivilegeContext.current()
    if inherited is not null:
        self.effectiveProfile = inherited          // invariant: == parent
    else:
        self.effectiveProfile = Profile.empty()    // daemon/headless path
    // invariant holds: prof(child) <= prof(parent)
\end{verbatim}

\paragraph{Why the invariant holds.}
The argument rests on three structural properties of the
implementation:

\begin{itemize}
    \item \textbf{Single assignment.} The child's
    \texttt{effectiveProfile} is set exactly once, in the
    constructor, from \texttt{SpawnPrivilegeContext.current()}.
    There is no subsequent code path that can replace it with a
    wider profile.

    \item \textbf{Try/finally guarantee.} The parent's
    \texttt{set}/\texttt{clear} pair is wrapped in
    \texttt{try}/\texttt{finally}, so the context is cleared
    even if the tool throws. This prevents thread-pool reuse
    from leaking a stale profile to an unrelated agent on the
    next turn---the only path that could violate the invariant
    across turns.

    \item \textbf{Inheritance by construction.}
    \texttt{InheritableThreadLocal} propagates the parent's
    value to the child virtual thread at construction time; the
    child cannot opt out, and no API exists for the child to
    \emph{replace} the inherited value with a wider one. The
    child can only \texttt{set} a narrower profile for its own
    sub-spawns, preserving the invariant down the spawn tree.
\end{itemize}

The combination of these properties ensures that the invariant
holds for every reachable state of the system, including the
exceptional and concurrent cases that an adversary might
attempt to exploit. A full formal verification in the style of
seL4~\cite{klein2009sel4} is left as future work; the argument
here is sufficient to show that no engineering code path
violates the property, which is the standard expected of
runtime-enforced invariants that are not formally verified.
